\begin{tikzpicture}[
  box/.style={draw=blue,fill=blue!5,rounded corners=2pt,text width=2.45cm,minimum height=.9cm,align=center,font=\footnotesize,inner sep=3pt},
  flow/.style={-{Stealth[length=4pt]},line width=.7pt,draw=black!60},
  label/.style={font=\scriptsize,align=center,fill=white,inner sep=1pt,text=black!60}]
\node[box] (g) at (1.5,0) {\textbf{Generator 团队}\\候选方案与测试};
\node[box,draw=forestgreen,fill=forestgreen!6] (checks) at (4.9,0) {\textbf{Orchestrator}\\数值门控};
\node[box,draw=darkred,fill=red!4] (c) at (8.3,0) {\textbf{Critic 团队}\\独立分维度审查};
\node[box,draw=forestgreen,fill=forestgreen!6] (o) at (11.7,0) {\textbf{Orchestrator}\\合并台账与路由};
\draw[flow] (g)--(checks);\draw[flow] (checks)--(c);\draw[flow] (c)--(o);
\draw[flow,dashed,draw=darkred] (checks.south) to[bend left=25] node[below,label,text=darkred]{有界修复}(g.south);
\node[box,text width=2.8cm,minimum height=1.3cm] (packet) at (1.5,-2.35) {\textbf{下一轮任务包}\\修复或再次确认};
\node[box,draw=forestgreen,fill=forestgreen!6,text width=3.1cm,minimum height=1.3cm] (accepted) at (6.6,-2.35) {\textbf{ACCEPTED}\\同一候选连续两轮通过};
\node[box,draw=darkred,fill=red!4,text width=3.0cm,minimum height=1.3cm] (stop) at (11.7,-2.35) {\textbf{未通过则停止}\\拒绝／耗尽\\停滞／受阻};
\draw[black!60,line width=.7pt] (o.south)--(11.7,-1.3)--(1.5,-1.3);
\draw[flow] (1.5,-1.3)--node[left,label]{继续}(packet.north);
\draw[flow] (6.6,-1.3)--node[right,label]{满足验收}(accepted.north);
\draw[flow] (11.7,-1.3)--node[right,label]{结束条件}(stop.north);
\draw[flow] (packet.west)--(-.25,-2.35)--(-.25,0)--(g.west);
\node[font=\scriptsize,align=center,text=black!60] at (6.6,-3.35) {修复超限记为 BLOCKED；更换候选方案重置计数；预算耗尽绝不视为通过。};
\end{tikzpicture}
